
\section{\BookRecordIndex{T0020}Frobenius 流形}\label{sec:11-5}

回到第~\ref{sec:11-2} 节的讨论，并仍假设 $(M,\omega)$ 是\BookRecordIndex{T0058}半正流形。为简单起见，我们只
考察 $M$ 的偶数次上同调，从而可以在流形而非\BookRecordIndex{T0075}超流形上工作。因此现在假设
$e_0:=1,e_1,\ldots,e_N$ 是 $H^{\mathrm{ev}}(M)$ 的一个由纯次数元素组成的基，并且
$e_1,\ldots,e_m$ 张成 $H^2(M)$。为方便起见，还假设
$e_N=\operatorname{PD}(\mathrm{pt})$ 是 $H^{2n}(M)$ 的正生成元。对\BookRecordIndex{T0274}旗流形和 Grassmann
流形等许多应用而言，只取偶数次类已经足够。

因为这里只想展现这一课题的概貌，我们再作一个简化假设：对所有
$a,b,c\in H^*(M)$，级数
\begin{equation}\label{eq:11-5-1}
  \sum_{A\in K^{\mathrm{eff}}}
  \operatorname{GW}^M_{A,3}(a,b,c)q^{\omega(A)}
\end{equation}
对绝对值充分小的 $q\in\mathbb C$ 收敛。这时可以像例~\ref{ex:11-1-4}(viii) 那样，取\BookRecordIndex{T0231}量子
系数环为 $\Lambda:=\mathbb C$。记
\[
  \mathcal H=\operatorname{QH}^{\mathrm{ev}}(M;\mathbb C)
  =\bigoplus_{i=0}^n\operatorname{QH}^{2i}(M;\mathbb C),
\]
即 $M$ 的\BookRecordIndex{T0229}量子上同调的偶数次部分。上述假设蕴含\BookRecordIndex{T0228}量子杯积在 $\mathcal H$ 上良定义。下面
说明 $\mathcal H$ 是如下意义的 \BookRecordIndex{T0019}Frobenius 代数。

\begin{definition}\label{def:11-5-1}
域 $\mathbb F=\mathbb R$ 或 $\mathbb C$ 上的一个 \emph{\BookRecordIndex{T0019}Frobenius 代数} $\mathcal A$，是
$\mathbb F$ 上带单位元 $1$ 的交换结合代数，并配备一个 $\mathbb F$-线性函数
$\alpha:\mathcal A\to\mathbb F$，使配对
\[
  \langle a,b\rangle:=\alpha(ab)
\]
非退化，也就是说，若对所有 $b$ 都有 $\langle a,b\rangle=0$，则 $a=0$。
\end{definition}

这个定义也可表述为：要求存在一个非退化双线性配对
$(a,b)\mapsto\langle a,b\rangle\in\mathbb F$，满足
\[
  \langle ab,c\rangle=\langle a,bc\rangle,
  \qquad a,b,c\in\mathcal A,
\]
然后令 $\alpha(a):=\langle1,a\rangle$。

对 $\mathcal H$，把乘法定义为\BookRecordIndex{T0228}量子杯积，并把 $\alpha$ 定义为在 $M$ 的基本同调类上积分。
于是对 $a_t=\sum_i t_i e_i$，有
\[
  \alpha(a_t):=\int_M a_t=t_N\in\mathbb C.
\]
相应配对与命题~\ref{prop:11-1-11}(iii) 中的配对不同，因为后者取值于底环 $R$，且不在
$\Lambda:=\mathbb C$ 上线性。不过，命题~\ref{prop:11-1-11} 的证明仍表明，上述定义确实
使 $\mathcal H$ 成为 $\mathbb C$ 上的有限维 \BookRecordIndex{T0019}Frobenius 代数。如果系数必须取 Novikov 环
而非域，就应适当修改上述概念；见 Abrams \srcite{7}。

\begin{remark}\label{rem:11-5-2}
$\operatorname{QH}^*(M)$ 的 \BookRecordIndex{T0019}Frobenius 代数结构有助于理解它的性质；可比较 Abrams
\srcite{7} 以及文献 \srcite{273} 附录中的讨论。
\end{remark}

依照 Dubrovin \srcite{93}，现在把\BookRecordIndex{T0228}量子杯积
\[
  \mathcal H\times\mathcal H\longrightarrow\mathcal H,
  \qquad (a,b)\longmapsto a*b
\]
解释成 $\mathcal H$ 的切丛上的联络，定义为
\begin{equation}\label{eq:11-5-2}
  \nabla_YX(a)=dX(a)Y(a)+X(a)*Y(a),
\end{equation}
其中 $a\in\mathcal H$，而 $X,Y:\mathcal H\to\mathcal H$ 是两个向量场。这里以显然的方式
把 $\mathcal H$ 的所有切空间与 $\mathcal H$ 等同，故向量场 $X$ 是映射
$a\mapsto X(a)\in T_a\mathcal H\equiv\mathcal H$，而它在 $a$ 处的导数
$dX(a):T_a\mathcal H\to T_{X(a)}\mathcal H$ 是 $\mathcal H$ 到自身的线性映射。下一个
引理说明，量子乘法的交换律可解释为\BookRecordIndex{T0248}挠率消失，而结合律可解释为曲率消失。

\begin{lemma}\label{lem:11-5-3}
\begin{enumerate}[label=(\roman*)]
\item 联络~\eqref{eq:11-5-2} 无挠，即
\[
  [X,Y]=\nabla_YX-\nabla_XY.
\]

\item 联络~\eqref{eq:11-5-2} 平坦，即
\[
  \nabla_X\nabla_YZ-\nabla_Y\nabla_XZ+\nabla_{[X,Y]}Z=0.
\]

\item 联络~\eqref{eq:11-5-2} 满足
\[
  \nabla_X\mathbf 1=X,
\]
其中 $\mathbf 1$ 表示由单位元 $1\in\mathcal H$ 定义的常值向量场。
\end{enumerate}
\end{lemma}

\begin{proof}
这里对 Lie 括号采用符号约定
\[
  [X,Y]=dX\cdot Y-dY\cdot X,
\]
其中 $dX(a)$ 表示映射 $X:\mathcal H\to\mathcal H$ 的微分，应把它看成
$\mathcal H$ 的线性变换，把 $Y$ 送到 $dX\cdot Y$。第一条立刻可见。为证明平坦性，注意
\begin{align*}
  \nabla_X\nabla_YZ
  &=\nabla_X(dZ\cdot Y+Z*Y)\\
  &=d(dZ\cdot Y+Z*Y)\cdot X+(dZ\cdot Y+Z*Y)*X\\
  &=d^2Z(Y,X)+dZ\cdot dY\cdot X+Z*(dY\cdot X)\\
  &\quad +(dZ\cdot X)*Y+(dZ\cdot Y)*X+(Z*Y)*X\\
  &=\nabla_{dY\cdot X}Z+d^2Z(X,Y)+Z*(Y*X)\\
  &\quad +(dZ\cdot X)*Y+(dZ\cdot Y)*X.
\end{align*}
第三条也是显然的。
\end{proof}

上述联络相当特殊，因为一形式
\[
  A:T\mathcal H\longrightarrow\operatorname{End}(\mathcal H),
  \qquad A_a(x)y=x*y
\]
是常值的，不依赖基点 $a$。一般的联络一形式
$A\in\Omega^1(\mathcal H,\operatorname{End}(\mathcal H))$ 可解释成一族乘法
\begin{equation}\label{eq:11-5-3}
  \mathcal H\times\mathcal H\longrightarrow\mathcal H,
  \qquad (x,y)\longmapsto x*_a y,
\end{equation}
它以 $\mathcal H$ 自身的元素为参数，所用公式是
\begin{equation}\label{eq:11-5-4}
  A_a(x)y=x*_a y,
\end{equation}
其中 $a\in\mathcal H$，且 $x,y\in T_a\mathcal H=\mathcal H$。把相应联络看成微分算子
\[
  C^\infty(\mathcal H,\mathcal H)
  \longrightarrow\Omega^1(\mathcal H,\mathcal H),
\]
它由 $\nabla=d+A$ 给出；更明确地说，对两个向量场
$X,Y:\mathcal H\to\mathcal H$，有
\begin{equation}\label{eq:11-5-5}
  \nabla_YX(a)=dX(a)Y(a)+X(a)*_aY(a).
\end{equation}
\BookRecordIndex{T0015}Dubrovin 联络~\eqref{eq:11-5-5} 的性质与乘法~\eqref{eq:11-5-3} 有如下关系。这里假设映射
\[
  \mathcal H\times\mathcal H\times\mathcal H\longrightarrow\mathcal H,
  \qquad (a,x,y)\longmapsto x*_a y
\]
是全纯的，并且关于 $x,y$ 复双线性。还要考察如下情形：这些乘法连同常值内积
$\langle\,\cdot,\cdot\,\rangle$（假设它是双线性而非埃尔米特型的），在每个切空间
$T_a\mathcal H=\mathcal H$ 上确定一个 \BookRecordIndex{T0019}Frobenius 代数结构。

\begin{lemma}\label{lem:11-5-4}
假设每个乘法 $(x,y)\mapsto x*_a y$ 都满足 Frobenius 条件。则：
\begin{enumerate}[label=(\roman*)]
\item 联络~\eqref{eq:11-5-5} 无挠，当且仅当乘法~\eqref{eq:11-5-3} 交换，即对所有
$a,x,y\in\mathcal H$，都有
\[
  x*_a y=y*_a x.
\]

\item 假设联络~\eqref{eq:11-5-5} 无挠。那么一形式~\eqref{eq:11-5-4} 是闭的，当且仅当
存在全纯函数 $S:\mathcal H\to\mathbb C$，使
\begin{equation}\label{eq:11-5-6}
  \langle x*_a y,z\rangle=\partial^3S_a(x,y,z).
\end{equation}

\item 假设联络~\eqref{eq:11-5-5} 无挠。那么一形式~\eqref{eq:11-5-4} 满足
$A\wedge A=0$，当且仅当所有乘法~\eqref{eq:11-5-3} 都结合。

\item 联络~\eqref{eq:11-5-5} 满足 $\nabla_X\mathbf 1=X$，当且仅当 $\mathbf 1$ 是所有
乘法~\eqref{eq:11-5-3} 的单位元。
\end{enumerate}
\end{lemma}

\begin{proof}
第 (i)、(iii)、(iv) 条直接由引理~\ref{lem:11-5-3} 得到。为证明第 (ii) 条，注意函数
\[
  \phi_a(x,y,z)=\langle x*_a y,z\rangle
\]
关于 $x,y,z$ 对称。此外，直接计算表明，它的导数
\[
  \psi_a(w,x,y,z)
  =\left.\frac{d}{dt}\right|_{t=0}\phi_{a+tw}(x,y,z)
\]
关于 $w,x,y,z$ 对称，当且仅当式~\eqref{eq:11-5-4} 给出的一形式 $A$ 是闭的。而
$\psi_a$ 的对称性又等价于存在全纯函数 $S:\mathcal H\to\mathbb C$，使
\[
  \phi_a(x,y,z)=\partial^3S_a(x,y,z).
\]
事实上，$S$ 的一个显式公式是
\[
  S(a)=\frac12\int_0^1(1-t)^2\phi_{ta}(a,a,a)\,dt.
\]
引理得证。
\end{proof}
\begin{remark}\label{rem:11-5-5}
前一引理第 (ii) 条的证明，可以用上同调
$\mathcal H=H^{\mathrm{ev}}(M;\mathbb C)$ 的复基 $e_0,\ldots,e_N$ 更明确地表述。
\sourcefootnote{原书在本注中把本节开头固定的基 $e_0,\ldots,e_N$ 的末指标三次改印成
$n$，从而写成 $\mathbb C^{n+1}$ 及求和至 $n$；但本章的 $n=\dim_{\mathbb R}M/2$，一般
并不等于 $\dim_{\mathbb C}\mathcal H-1$，故此处统一恢复为 $N$。}
于是可用同构
\[
  \mathbb C^{N+1}\longrightarrow\mathcal H,
  \qquad x\longmapsto\sum_i x^i e_i
\]
把 $\mathcal H$ 与 $\mathbb C^{N+1}$ 等同。定义函数
$A_{ijk}:\mathbb C^{N+1}\to\mathbb C$ 为
\[
  A_{ijk}(x)=\langle e_i*_a e_j,e_k\rangle,
  \qquad a=\sum_{i=0}^N x^i e_i.
\]
按假设，这些函数是全纯的。它们在 $i,j,k$ 的置换下对称，当且仅当乘法~\eqref{eq:11-5-3}
交换并满足 Frobenius 条件。若此条件成立，则存在全纯函数
$S:\mathbb C^{N+1}\to\mathbb C$，满足
\[
  \frac{\partial^3S}{\partial x^i\partial x^j\partial x^k}=A_{ijk},
\]
当且仅当导数 $\partial_\ell A_{ijk}$ 在 $i,j,k,\ell$ 的置换下对称。$S$ 的一个显式公式是
\[
  S(x)=\frac12\int_0^1(1-t)^2
  \sum_{i,j,k}A_{ijk}(tx)x^ix^jx^k\,dt.
\]
\end{remark}

若全纯函数
\[
  S:\mathcal H\longrightarrow\mathbb C
\]
的三阶导数满足式~\eqref{eq:11-5-6}，就称它为 \BookRecordIndex{T0015}Dubrovin 联络~\eqref{eq:11-5-5} 的
\emph{\BookRecordIndex{T0301}势函数}。由引理~\ref{lem:11-5-4}，这样的函数存在，当且仅当该联络满足 Frobenius
条件、无挠并且联络一形式是闭的。反过来，给定\BookRecordIndex{T0301}势函数
$S:\mathcal H\to\mathbb C$，并用式~\eqref{eq:11-5-6} 定义乘法~\eqref{eq:11-5-3}，则
这些乘法自动交换并满足 Frobenius 条件。此外，$\mathbf1$ 是单位元，当且仅当 $S$ 满足
\begin{equation}\label{eq:11-5-7}
  \partial^3S_a(\mathbf1,x,y)=\langle x,y\rangle.
\end{equation}
结合律转化为 WDVV 方程
\begin{equation}\label{eq:11-5-8}
  \sum_{i,j}\partial^3S_a(w,x,e_i)g^{ij}
  \partial^3S_a(e_j,y,z)
  =
  \sum_{i,j}\partial^3S_a(w,z,e_i)g^{ij}
  \partial^3S_a(e_j,x,y),
\end{equation}
其中 $e_i$ 表示 $\mathcal H$ 的一个基，$g_{ij}:=\langle e_i,e_j\rangle$，而
$g^{ij}$ 表示逆矩阵。

\begin{remark}\label{rem:11-5-6}
像注~\ref{rem:11-5-5} 那样选取 $\mathcal H$ 的基，并令 $e_0=1$。则
条件~\eqref{eq:11-5-7} 具有形式
\begin{equation}\label{eq:11-5-9}
  \frac{\partial^3S}{\partial x^0\partial x^i\partial x^j}=g_{ij},
\end{equation}
而 WDVV 方程可写成
\[
  \sum_{\nu,\mu}
  \frac{\partial^3S}{\partial x^i\partial x^j\partial x^\nu}
  g^{\nu\mu}
  \frac{\partial^3S}{\partial x^\mu\partial x^k\partial x^\ell}
  =
  \sum_{\nu,\mu}
  \frac{\partial^3S}{\partial x^i\partial x^\ell\partial x^\nu}
  g^{\nu\mu}
  \frac{\partial^3S}{\partial x^\mu\partial x^j\partial x^k}
\]
对所有 $i,j,k,\ell$ 成立。
\end{remark}

具有这种结构的流形称为 \emph{\BookRecordIndex{T0020}Frobenius 流形}。更确切地说，一个（复）\BookRecordIndex{T0020}Frobenius
流形是三元组 $(\mathcal H,g,\nabla)$，由一个光滑复流形、一个平坦复双线性度量
$g(x,y):=\langle x,y\rangle$ 和第二个无挠复线性联络 $\nabla$（\BookRecordIndex{T0015}Dubrovin 联络）组成；它与
度量的 \BookRecordIndex{T0034}Levi--Civita 联络 $d$ 由
\[
  \nabla=d+A
\]
联系，其中 $A$ 是闭的且 $A\wedge A=0$。于是 $\mathcal H$ 的切空间都是 \BookRecordIndex{T0019}Frobenius
代数，并且联络具有\BookRecordIndex{T0301}势函数 $S$。下一个引理列出这些条件的不同等价表述。

\begin{lemma}\label{lem:11-5-7}
设 $S:\mathcal H\to\mathbb C$ 是满足~\eqref{eq:11-5-7} 的全纯函数，并用
式~\eqref{eq:11-5-3} 为 $a\in\mathcal H$ 定义量子乘法 $x*_a y$。则下列条件等价：
\begin{enumerate}[label=(\roman*)]
\item 乘法 $x*_a y$ 结合；
\item \BookRecordIndex{T0015}Dubrovin 联络~\eqref{eq:11-5-5} 平坦；
\item \BookRecordIndex{T0301}势函数 $S$ 满足 WDVV 方程~\eqref{eq:11-5-8}。
\end{enumerate}
\end{lemma}

\begin{proof}
第 (i) 与 (ii) 条等价，因为由引理~\ref{lem:11-5-4}(ii)，联络一形式
$A$ 是闭的；而由该引理第 (iii) 条，条件 $A\wedge A=0$ 等价于结合律。由于 $\nabla$ 的
曲率是取值于自同态的二形式
\[
  F_A=dA+A\wedge A,
\]
事实上还得到：结合律等价于所有联络
\[
  \nabla_\lambda=d+\lambda A,
  \qquad \lambda\in\mathbb R,
\]
都平坦。

为证明第 (i) 与 (iii) 条等价，注意
\[
  x=\sum_{i,j}\langle x,e_i\rangle g^{ij}e_j,
  \qquad
  \langle x,y\rangle
  =\sum_{i,j}\langle x,e_i\rangle g^{ij}\langle e_j,y\rangle
\]
对所有 $x,y\in\mathcal H$ 成立。这蕴含
\begin{align*}
  &\sum_{i,j}\partial^3S_a(w,x,e_i)g^{ij}
  \partial^3S_a(e_j,y,z)\\
  &\qquad
  =\sum_{i,j}\langle w*_a x,e_i\rangle g^{ij}
    \langle e_j,y*_a z\rangle\\
  &\qquad
  =\langle w*_a x,y*_a z\rangle\\
  &\qquad
  =\langle w,x*_a(y*_a z)\rangle.
\end{align*}
\sourcefootnote{原书第二行把逆度量展开式中的基向量误印为未定义的 $g_j$；这里唯一应为
$e_j$。}
所以结合律等价于式~\eqref{eq:11-5-8} 左端在 $x,y,z$ 的置换下对称。引理得证。
\end{proof}

如果考察不依赖 $a$ 的杯积量子形变（即\BookRecordIndex{T0349}小量子乘法），相应\BookRecordIndex{T0301}势函数 $S$ 就是三次多项式。
由\BookRecordIndex{T0228}量子杯积的定义，这个三次多项式为
\[
  S_3(a)=\frac1{3!}\sum_{A\in K^{\mathrm{eff}}}
  \operatorname{GW}^M_{A,3}(a,a,a).
\]
（为使该式有意义，我们假设例~\ref{ex:11-1-4}(viii) 中的级数收敛。）

\emph{\BookRecordIndex{T0083}大量子上同调}的\BookRecordIndex{T0301}势函数，是完整的\BookRecordIndex{T0144}格罗莫夫--威滕势
\[
  \Phi(q,t):=\sum_\alpha\sum_d\frac{\varepsilon(\alpha)}{\alpha!}
  \operatorname{GW}^M_{A_d,|\alpha|}(e^\alpha)q^dt^\alpha.
\]
定理~\ref{thm:11-2-4} 已说明，$\Phi$ 是变量
\[
  (q,t)\in\mathbb C^m\times\mathcal H=\mathbb C^{m+N+1}
\]
的形式幂级数环中 WDVV 方程~\eqref{eq:11-5-8} 的解。
\sourcefootnote{原书把 $\mathbb C^m\times\mathcal H$ 的维数印作 $m+N$；但本节固定
$\mathcal H$ 的基为 $e_0,\ldots,e_N$，故其复维数为 $N+1$，总维数唯一应为
$m+N+1$。}
如果 $\Phi(q,t)$ 在 $(0,0)\in\mathbb C^m\times\mathcal H$ 的某个开邻域中收敛，那么适当
选取 $q\in\mathbb C^m$ 后，它就在原点的某个邻域
$\mathcal U\subset\mathcal H$ 上定义全纯函数 $S_\Phi$。此外，\BookRecordIndex{T0345}弦方程正是
式~\eqref{eq:11-5-9}。因此 $\mathcal U$ 成为 \BookRecordIndex{T0020}Frobenius 流形。$\mathcal H$ 上相应的乘法族
$x*_a y$（$a\in\mathcal U$）称为\BookRecordIndex{T0083}大量子上同调。明确地说，$x*_a y$ 由
\[
  \langle x*_a y,z\rangle
  :=\sum_p\frac1{p!}\sum_A
  \operatorname{GW}^M_{A,p+3}(x,y,z,a,\ldots,a)
\]
定义。如果这个收敛假设不成立，就必须把 $\mathcal U$ 看成 $\mathcal H$ 中原点的形式邻域；
见 Manin \srcite{286, \S5.2.1}。

一般说来，\BookRecordIndex{T0083}大量子上同调比小\BookRecordIndex{T0229}量子上同调包含多得多的信息。不过，Kontsevich 和 Manin
在文献 \srcite{216} 中证明：若 $M$ 的上同调环 $H^*(M)$ 由 $H^2(M)$ 生成，则大量子
乘法不含新信息。

\begin{theorem}[Kontsevich--Manin]\label{thm:11-5-8}
设 $(M,\omega)$ 是\BookRecordIndex{T0058}半正\BookRecordIndex{T0358}辛流形，并且 $H^*(M)$ 作为环由 $H^2(M)$ 生成。那么
$\operatorname{QH}^*(M;\Lambda)$ 的结构常数决定\BookRecordIndex{T0144}格罗莫夫--威滕势。
\end{theorem}

对一般流形，只能说\BookRecordIndex{T0083}大量子上同调在 $H^2(M;\mathbb C)$ 上的限制由小数\BookRecordIndex{T0229}量子上同调决定。
换言之，$\mathcal H$ 上由 $a\in H^2(M;\mathbb C)$ 参数化的乘法族 $*_a$，由小乘法 $*$
决定。这等价于说，\BookRecordIndex{T0143}格罗莫夫--威滕不变量
\[
  \operatorname{GW}^M_{A,p+3}(x,y,z,a,\ldots,a),
  \qquad a\in H^2(M;\mathbb C),
\]
由不变量 $\operatorname{GW}^M_{A,3}(x,y,z)$ 决定；这是\BookRecordIndex{T0077}除子方程的直接推论。

\begin{exercise}\label{ex:11-5-9}
求乘法族 $h*_t h$ 在
\[
  \mathcal H:=\operatorname{QH}^2(Z_5;\mathbb C)
\]
上的显式公式，其中 $Z_5$ 是第~\ref{subsec:11-3-4} 小节讨论的
$\mathbb{CP}^4$ 中五次超曲面，而 $t$ 是一维空间 $\mathcal H$ 上的坐标。
\end{exercise}

\BookRecordIndex{T0020}Frobenius 流形有丰富而有趣的结构。它们出现在孤立奇点的万有形变、可积系统，以及更一般地，
带有自然平坦联络的\BookRecordIndex{T0246}模空间的研究中。这里无法描述一般理论，只会指出当前模型的若干特殊性质。

先来描述与\BookRecordIndex{T0077}除子方程对应的附加结构。\BookRecordIndex{T0077}除子方程蕴含，$\Phi(q,t)$ 是如下两项之和：一项是变量
\[
  t_0,\quad q_i e^{t_i}\ (i=1,\ldots,m),\quad
  t_\nu\ (\nu>m)
\]
的函数，另一项是变量 $t_\nu$（$0\le\nu\le N$）的二次多项式。因此，在参数化
$H^2(M;\mathbb C)$ 的变量 $t_1,\ldots,t_m$ 上加 $2\pi i$ 的整数倍，不会改变
$\partial^3S_\Phi$，从而也不改变联络。故 $\mathcal U\subset\mathcal H$ 上的 \BookRecordIndex{T0020}Frobenius
流形结构下降到商空间
\[
  \overline{\mathcal H}
  :=\frac{\operatorname{QH}^{\mathrm{ev}}(M;\mathbb C)}
  {H^2(M;2\pi i\mathbb Z)}
  =\mathbb C e_0\oplus H^2(M;\mathbb C/2\pi i\mathbb Z)
  \oplus\bigoplus_{i>1}\operatorname{QH}^{2i}(M;\mathbb C)
\]
的一个子集上。另一个推论是，$\mathcal U$ 上的 \BookRecordIndex{T0020}Frobenius 流形结构并不本质依赖为变量
$q_i$ 所选的值 $q_i\in\mathbb C$，因为改变这些值可由重标度 $t_i$（$1\le i\le m$）来
补偿。反过来，改变 $t_i$ 的值也可通过改变 $q_i\in\mathbb C$，等价地通过改变\BookRecordIndex{T0363}辛形式的类，
来补偿。
\footnote{这里要提醒读者注意记号：尽管我们的变量称为 $q,t$，它们与系数环
$\Lambda^{\mathrm{univ}}[q,q^{-1}]$ 中的变量 $q,t$ 含义不同。特别地，由式~\eqref{eq:11-5-1}
可见，现在的值 $q\in\mathbb C$ 携带了\BookRecordIndex{T0363}辛形式的信息。}
所以，\BookRecordIndex{T0019}Frobenius 代数族
$T_a\mathcal H$（$a\in H^2(M;\mathbb C)$）可与一族适当\BookRecordIndex{T0363}辛形式（或 K\"ahler 形式）
$\omega_q$ 所对应的 $(M,\omega_q)$ 的小\BookRecordIndex{T0229}量子上同调 \BookRecordIndex{T0019}Frobenius 代数族等同；可比较
习题~\ref{ex:11-5-9}。

把这一说明应用于 Calabi--Yau 三维流形。第~\ref{subsec:11-3-4} 小节已说明，
$H^2(M;\mathbb C)$ 上有一个依赖于\BookRecordIndex{T0363}辛形式类 $[\omega_q]$ 的诱导量子乘法结构。像上面那样
使用\BookRecordIndex{T0077}除子方程，可以说明该结构下降到 $M$ 的复化 \BookRecordIndex{T0030}K\"ahler 锥
\[
  \mathcal K:=H^2(M;\mathbb C/2\pi i\mathbb Z),
\]
并在平凡丛 $\mathcal K\times\mathcal H\to\mathcal K$ 上定义 \BookRecordIndex{T0015}Dubrovin 联络的一个类似物。
这个联络记录当类 $[\omega_q]$ 在 $\mathcal K$ 中变化时\BookRecordIndex{T0349}小量子乘法的变化，称为
\emph{\BookRecordIndex{T0010}A 模型联络}。虽然它是平坦的，却绕奇点 $q=0$ 具有非平凡\BookRecordIndex{T0095}单值化。这个联络恰与 Hodge
结构形变的 Gauss--Manin 联络互为镜像；见 Cox--Katz \srcite{76}。

\BookRecordIndex{T0144}格罗莫夫--威滕势 $\Phi$ 的另一个重要事实是：当
$\deg(t_\nu):=2-\deg(e_\nu)$ 时，它是次数 $6-\dim M$ 的齐次函数。此外，我们在
习题~\ref{ex:11-2-6} 中证明，存在微分算子 $\mathcal E$，使
\[
  \mathcal E\Phi=(3-n)\Phi.
\]
$\mathcal E$ 的公式包含诸项 $\sum c^iq_i\partial/\partial q_i$（其中
$\sum c^i e_i$ 是第一陈类），所以 $\mathcal E$ 不是 $\mathcal H$ 上的向量场。但是，把
这些项换成 $\sum_i c^i\partial/\partial t_i$，就得到 $\mathcal H$ 上的向量场
\[
  E(t)=\sum_{i=1}^m c^i e_i
  +\sum_{\nu=0}^N\left(1-\frac{\deg(e_\nu)}2\right)t_\nu e_\nu.
\]
（仍像前面那样，把向量场看成映射 $\mathcal H\to\mathcal H$；因此 $e_\nu$ 对应微分算子
$\partial/\partial e_\nu$。）\BookRecordIndex{T0077}除子方程蕴含 $E\Phi-(3-n)\Phi$ 是二次多项式。因此 $E$ 在
Frobenius 结构上作共形作用；见下面的习题~\ref{ex:11-5-10}。当 $E$ 满足称为\BookRecordIndex{T0057}半单性的
附加条件时，可以在 \BookRecordIndex{T0020}Frobenius 流形 $\mathcal H$ 上引入新的典范坐标。这是\BookRecordIndex{T0229}量子上同调所产生
的 \BookRecordIndex{T0020}Frobenius 流形结构中的重要部分。更多信息见 Manin \srcite{286}。

\begin{exercise}\label{ex:11-5-10}
设 $E$ 如上定义。验证对 $\mathcal H$ 上所有向量场 $X,Y$，都有
\[
  E\bigl(g(X,Y)\bigr)+g([E,X],Y)+g(X,[E,Y])=Dg(X,Y),
\]
其中 $D:=2-\dim M/2$，并且
\[
  -[E,X*Y]+[E,X]*Y+X*[E,Y]=X*Y.
\]
第一式说明 $E$ 是共形的，即 $\mathcal L_E(g)=Dg$。这里 $\mathcal L_E$ 表示 Lie 导数，
依照本书对向量场的符号约定定义为
\[
  \mathcal L_E(X):=-[E,X].
\]
第二式说明 $E$ 在乘法结构上也作共形作用，即
$\mathcal L_E(*)=d_0(*)$，其中 $d_0=1$。满足这两个恒等式的向量场称为 \BookRecordIndex{T0016}Euler 场。
\sourcefootnote{原书把 Lie 导数误印为 $\mathcal L_E(X):=-[X,Y]$，并把第二个恒等式首项
印成 $+[E,X*Y]$。按本书约定 $[X,Y]=dX\cdot Y-dY\cdot X$，前者唯一应为
$-[E,X]$；将这一导数作用于双线性乘法，后者首项也唯一应带负号。}
\end{exercise}

最后作几点说明，指出 Givental \srcite{150} 关于环面完全交镜像猜想的工作，以及他与
Coates \srcite{71,72,152} 进一步发展的 Frobenius 结构拉格朗日锥图景所深入研究的方向。

首先回忆，在引理~\ref{lem:11-5-7} 的证明中，我们注意到满足 WDVV 方程的\BookRecordIndex{T0301}势函数 $S$ 定义
一族平坦的 \BookRecordIndex{T0015}Dubrovin 联络
\[
  \nabla_\lambda:=d+\lambda A,
  \qquad \lambda\in\mathbb C.
\]
可以把 $\lambda$ 看成环 $\mathbb C[\lambda,\lambda^{-1}]$ 的元素，而该环又可解释成
$\mathbb C\setminus\{0\}$ 上的函数环（Manin 的观点），或 $BS^1$ 的局部化上同调环
（Givental 的观点）。在某种意义下，这在代数上相当于用 \BookRecordIndex{T0003}$J$-全纯曲线
$u:S^2\to M$ 的图像
\[
  \widetilde u:S^2\longrightarrow S^2\times M
\]
代替 $u$，然后使用关于 $S^2\times M$ 上圆周作用的等变理论，其中这个作用旋转球面。不过，
在文献 \srcite{146} 中，Givental 采用 $M$ 的环路空间上的等变 弗洛尔理论来思考。尽管这种
版本的 弗洛尔理论尚未建立，其基本思想已显示出强大力量，并启发了 Givental 随后在这一
领域的全部工作。

其次，Givental 发现，对所有 $\lambda$ 都 $\nabla_\lambda$-平坦的 $T\mathcal H$ 截面 $S$
在理论中扮演重要角色。他把这些平坦截面的适当分量组合成 \BookRecordIndex{T0004}$J$-函数；这是他研究镜像猜想的
关键成分，见 Cox--Katz \srcite{76}。理论的这两个方面彼此相关：Givental 找到了用带 $k$
个\BookRecordIndex{T0066}标记点的\BookRecordIndex{T0378}引力后裔表示平坦截面的公式，并在文献 \srcite{147} 中借助
$S^2\times M$ 上的 $S^1$-等变理论，通过局部化来解释它。（\BookRecordIndex{T0378}引力后裔来自不动点集\BookRecordIndex{T0117}法丛的
Euler 类。）

因此，考察\BookRecordIndex{T0168}后裔势 $\mathcal F$（亏格零\BookRecordIndex{T0378}引力后裔的生成函数），而不是这里考察的
\BookRecordIndex{T0144}格罗莫夫--威滕势 $\Phi$，会得到更令人满意的理论。作适当的线性变量替换（\BookRecordIndex{T0262}膨胀子平移）
之后，可把 $\mathcal F$ 看成空间 $\mathcal H[z]$ 上的（形式）函数，其中
$\mathcal H[z]$ 是以 $\mathcal H:=H^*(M;\mathbb C)$ 为系数的多项式函数空间。（这里应把
$z$ 解释成 $1/\lambda$，其中 $\lambda$ 参数化 \BookRecordIndex{T0015}Dubrovin 联络族。）它满足\BookRecordIndex{T0345}弦方程、\BookRecordIndex{T0077}除子方程
和 WDVV 方程的相应版本。在与 Coates 合作的工作中，Givental \srcite{71} 把这些基本方程
完全解释为余切丛 $T^*(\mathcal H[z])$ 中\BookRecordIndex{T0219}拉格朗日子流形 $\mathcal L$ 的几何性质，其中
$\mathcal L$ 是 $d\mathcal F$ 的图像。这为理解亏格零\BookRecordIndex{T0143}格罗莫夫--威滕不变量的结构提供了
非常有力的框架。特别地，由于 $\mathcal F$ 在所选坐标中是二次齐次的，$\mathcal L$ 是锥；
而 \BookRecordIndex{T0004}$J$-函数就是 $\mathcal L$ 与仿射子空间
\[
  \left.T^*(\mathcal H[z])\right|_{-z+\mathcal H}
\]
的交，后者由把余切丛限制到有限维切片 $-z+\mathcal H$ 得到。此外，证明环面完全交的
镜像定理时所需的变换，可以通过环境线性空间 $T^*(\mathcal H[z])$ 的辛变换及其对
$\mathcal L$ 的作用来理解。高亏格不变量的相应图景尚未得到理解。
